\documentclass[]{article}
\usepackage{amsmath}
\usepackage{amsfonts}
\usepackage{graphicx}
\usepackage{url}
\usepackage{listings}
\usepackage[bookmarks]{hyperref}
\usepackage[a4paper, total={7in, 8in}]{geometry}

\lstset{basicstyle=\footnotesize\ttfamily,breaklines=true}

\DeclareMathOperator{\sign}{sign}
\DeclareMathOperator{\Tr}{Tr}
\newcommand{\realevol}{\texttt{realevol}}
\newcommand{\bra}[1]{\ensuremath{\langle#1|}}
\newcommand{\ket}[1]{\ensuremath{|#1\rangle}}
\newcommand{\aver}[1]{\ensuremath{\langle#1\rangle}}

% Title Page
\title{Lanczos short time propagation solver: implementation notes}
\author{Igor Krivenko}

\begin{document}
\maketitle

The Lanczos short time propagation solver (\realevol) is designed to solve models of many-body quantum systems with interacting fermions and bosons out of thermodynamical equilibrium. It is based on TRIQS library version 1.4.x. Similarly to other TRIQS applications, \realevol\ is written in C++14 and comes with a Python interface, which allows to run physical simulations by writing simple Python scripts.

Besides TRIQS, \realevol\ has the following prerequisites:
\begin{itemize}
    \item Boost C++ libraries, version 1.53 or newer (headers only);
    \item A working MPI implementation;
    \item A recent version of ARPACK-NG (\url{https://github.com/opencollab/arpack-ng});
    \item \texttt{triqs\_arpack} version 0.5, coinstalled with TRIQS
    (\url{https://github.com/krivenko/triqs_arpack/tree/0.5}).
\end{itemize}

\realevol\ does runtime evaluation of mathematical expressions when computes
matrix elements of time-dependent Hamiltonians. A backend for this feature is
provided by the C++ Mathematical Expression Parsing And Evaluation Library,
a.k.a. ExprTk from Arash Partow (\url{https://github.com/ArashPartow/exprtk}).
A copy of ExprTk is bundled to the source codes.

In order to ease installation of \realevol\ and of its prerequisites, I wrote a
shell script called \texttt{build\_realevol\_and\_deps.sh} (to be found in \texttt{/scripts} subdirectory). When run, it tries to find the source codes of \realevol\ in subdirectory \texttt{\$(pwd)/realevol.git}. It also creates two more subdirectories in the current directory,
\begin{itemize}
    \item \texttt{\$(pwd)/build}, to place distributives and build files of TRIQS, ARPACK-NG, \texttt{triqs\_arpack} and \realevol;
    \item \texttt{\$(pwd)/installed}, as the installation directory for \realevol\ and the prerequisites.
\end{itemize}

The installation procedure can be tweaked by changing the header section of
\texttt{build\_realevol\_and\_deps.sh}:

\lstinputlisting[language=bash, firstline=3, lastline=11]
    {../../scripts/build_realevol_and_deps.sh}

\section{\label{sec:formulation}Problem formulation}
\subsection{\label{sec:formulation:space}States space}

Let us consider a quantum system of many identical particles. Its states space $\mathcal{H}$ (Fock space) is generated by fixed finite sets of fermionic creation operators $\{c^\dagger_1,\ldots,c^\dagger_{N_f}\}$ and bosonic creation operators $\{b^\dagger_1,\ldots,b^\dagger_{N_b}\}$. The total space $\mathcal{H}$ factorizes as
\begin{equation}
    \mathcal{H} = \mathcal{H}_f \otimes \mathcal{H}_b,
\end{equation}
where $\mathcal{H}_f$ is finite-dimensional fermionic states space, $\dim(\mathcal{H}_f) = 2^{N_f}$ (as follows from Pauli exclusion principle). The bosonic states space $\mathcal{H}_b$ is formally infinitely-dimensional, as the number of bosons created by every operator $b^\dagger_i$ is unbounded.
\realevol\ is an exact-diagonalization type solver, which is bound to work with vector spaces of finite dimensions. Therefore, one has to artificially truncate $\mathcal{H}_b$ so that its dimension is finite. Truncation is done by setting a cutoff on the number of bosonic excitations in each mode,
\begin{equation}
    b^\dagger_i\ket{n_i^\mathrm{max}} = 0, \quad i=1,\ldots,N_b.
\end{equation}
The truncation constants $n_i^\mathrm{max}$ are chosen by user as one of input parameters. For technical reasons, the user is asked to provide a list of positive numbers called \texttt{bits\_per\_boson} and defined such that
\begin{equation}
    n_i^\mathrm{max} = 2^\mathtt{bits\_per\_boson[{\it i}]} - 1.
\end{equation}

According to a standard TRIQS notation, the abstract index $i$ must take a form
of a pair \texttt{(block\_index,inner\_index)} in Python scripts. Each element
of the pair can either be an integer, or a string. The rationale behind this
notation is that Green's function containers share the same type of indices in TRIQS. In many physical problems Green's functions turn out to be block-diagonal. This is why there must be one index to address a diagonal block,
and another index to choose a column/row within that block. For example, local
dynamics of a 3-orbital quantum impurity without spin-orbital coupling can be
described by a GF with block indices \texttt{'up'}, \texttt{'down'} (spin
projections) and inner indices \texttt{1}, \texttt{2}, \texttt{3} (orbitals). Therefore, there are 6
index combinations used to label operators $c^\dagger$/$c$:
\texttt{('up',1)}, \texttt{('up',2)}, \texttt{('up',3)}, \texttt{('down',1)},
\texttt{('down',2)} and \texttt{('down',3)}.

\subsection{\label{sec:formulation:initial}Initial state}
At the initial time moment ($t=0$) the system is prepared in a known state characterized by density matrix $\hat\rho_0$. \realevol\ supports three
qualitatively different types of initial states.

\begin{enumerate}
    \item Pure state, $\hat\rho_0 = \ket{\psi}\bra{\psi}$.

    Instead of specifying $\ket{\psi}$ as a list of coefficients w.r.t. some basis in $\mathcal{H}$, user is expected to provide a generating operator $\hat P$,
    \begin{equation}
        \ket{\psi} = \frac{\hat P\ket{0}}
                    {\sqrt{\bra{0}\hat P^\dagger \hat P\ket{0}}}.
    \end{equation}
    (A runtime error will occur if $\hat P\ket{0} = 0$).
    This approach is more convenient in practice, because it does not require knowledge about the particular choice of Fock basis in $\mathcal{H}$.

    \item Ground state(s) of a given Hamiltonian $\hat H_0$.
    \begin{equation}
        \hat\rho_0 = \frac{1}{g_\mathrm{gs}}\sum_{\psi\in\{\mathrm{gs}\}} \ket{\psi}\bra{\psi}.
    \end{equation}
    $g_\mathrm{gs}$ stands for the degeneracy degree of the ground state.

    \item Thermal state with a given Hamiltonian $\hat H_0$ at temperature $T$.
    \begin{equation}
        \hat\rho_0 = \frac{\exp(-\hat H_0/T)}{Z},
        \quad Z = \Tr[\exp(-\hat H_0/T)].
    \end{equation}
\end{enumerate}

Initial state (2) can be seen as the zero temperature limit of (3).

\subsection{\label{sec:formulation:observables}Observables}

Evolution of the system at $t>0$ is governed by a time-dependent Hamiltonian $\hat H(t)$.

There are currently three time-dependent observables \realevol\ can compute \cite{NoneqDMFT}.
\begin{itemize}
    \item Greater Green's function of fermions,
    \begin{equation}
        G^>_{ij}(t,t') = -i\frac{1}{\hbar}\aver{c_i(t) c^\dagger_j(t')} =
        -i\frac{1}{\hbar}\Tr\left[\hat\rho_0 c_i(t) c^\dagger_j(t')\right].
    \end{equation}

    \item Lesser Green's function of fermions,
    \begin{equation}
        G^<_{ij}(t,t') = i\frac{1}{\hbar}\aver{c^\dagger_j(t') c_i(t)} =
        i\frac{1}{\hbar}\Tr\left[\hat\rho_0 c^\dagger_j(t') c_i(t)\right].
    \end{equation}

    \item Generalized susceptibility ($n_i = c^\dagger_i c_i$),
    \begin{equation}
        \chi_{ij}(t,t') = -i\frac{1}{\hbar}\aver{n_i(t) n_j(t')} =
        -i\frac{1}{\hbar}\Tr\left[\hat\rho_0 n_i(t) n_j(t')\right].
    \end{equation}
\end{itemize}

Here, $c_i(t)$, $c^\dagger_i(t)$ and $n_i(t)$ are Heisenberg operators related to their Schr\"odinger counterparts as
\begin{equation}
    \hat A(t) = \hat U(0,t)\hat A \hat U(t,0).
\end{equation}
Following a standard derivation, one can show that the time evolution operator $\hat U(t,t')$ is given by a time-(anti)ordered exponential,
\begin{equation}
    \hat U(t,t') = \left\{
        \begin{array}{ll}
        \mathbb{T}_t\exp\left[
        -\frac{i}{\hbar}\int_{t'}^t d\tilde t\hat H(\tilde t) \right], &t>t',\\
        \mathbb{\bar T}_t\exp\left[
        -\frac{i}{\hbar}\int_{t'}^t d\tilde t\hat H(\tilde t) \right], &t<t'.
        \end{array}
    \right.
\end{equation}

\textbf{Warning}: Even though the studied system can contain bosonic degrees of freedom, computation of correlators of $b^\dagger_i$/$b_i$ is not implemented.

\section{\label{sec:init_state}\texttt{InitState}: Representation of the initial state}

Any initial state $\hat\rho_0$ supported by \realevol\ can be represented as a finite sum of weighted pure states.
\begin{equation}
    \hat\rho_0 = \sum_m w_m \ket{\psi_m}\bra{\psi_m},\quad\sum_m w_m = 1.
\end{equation}
The way state vectors $\ket{\psi_m}$ and their respective weights $w_m$ are chosen, varies between the three cases listed in \ref{sec:formulation:initial} and will be discussed below. Before we proceed with this discussion, we need to introduce the notion of states space partition.

The full states space of the problem can, in general, be partitioned into a direct sum of subspaces,
\begin{equation}
    \mathcal{H} = \mathcal{H}_1 \oplus \mathcal{H}_2 \oplus \ldots \oplus
                  \mathcal{H}_S,\quad 1\le S\le \dim(\mathcal{H}).
\end{equation}
If a given state $\ket{\psi}\in\mathcal{H}_s$, it is possible to save memory by storing only expansion coefficients of $\ket{\psi}$ over the basis vectors of
$\mathcal{H}_s$. This optimization technique is implemented in \texttt{InitState}, the Python object used to store $\hat\rho_0$ in \realevol\
(the corresponding C++ type is called \texttt{init\_state}). An instance of
\texttt{InitState} contains information about space partition ($S$, and a list of basis vectors for each $\mathcal{H}_s$) and a list of all contributing state/weight pairs accompanied by indices of the respective subspaces.
\begin{equation}
    \mathtt{InitState} = \{\mathcal{H}_s\}, \{(\ket{\psi_m}, w_m, s)\}.
\end{equation}

\section{\label{sec:make_pure_init_state}Function \texttt{make\_pure\_init\_state()}}

Python function \texttt{make\_pure\_init\_state()} creates a pure initial state (see Sec.~\ref{sec:formulation:initial}, point 1). It takes four arguments.
\begin{itemize}
    \item A Python object of type \texttt{realevol.operators.Operator}
    representing the generating operator $\hat P$.

    $\hat P$ must be time-independent and such that $\hat P\ket{0} \neq 0$.
    \item Index set for fermions, i.e. a \texttt{set} of \texttt{(block\_index,inner\_index)} pairs to label all operators $c^\dagger_i$/$c_i$.

    \item Index set for bosons, i.e. a \texttt{set} of \texttt{(block\_index,inner\_index)} pairs to label all operators $b^\dagger_i$/$b_i$.

    \item \texttt{bits\_per\_boson}, number of bits per bosonic
    degree of freedom passed as Python dictionary\\
    \texttt{\{(block\_index,inner\_index):int\}}. Must be empty if
    system contains no bosons.

\end{itemize}

The two index sets and \texttt{bits\_per\_boson} together define the full Fock
space $\mathcal{H}$ (see Sec.~\ref{sec:formulation:space}).

\texttt{make\_pure\_init\_state()} returns an instance of \texttt{InitState}
containing one triplet
\begin{equation}
    \left(
        \frac{\hat P\ket{0}}{\sqrt{\bra{0}\hat P^\dagger \hat P\ket{0}}},
        1,
        1
    \right),
\end{equation}
and two subspaces, $\mathcal{H}=\mathcal{H}_1\otimes\mathcal{H}_2$.
$\mathcal{H}_1$ is spanned by all basis vectors of $\mathcal{H}$ that have non-zero overlaps with $\hat P\ket{0}$. The rest of the basis vectors span
$\mathcal{H}_2$.

\section{\label{sec:make_equilibrium_init_state}Function \texttt{make\_equilibrium\_init\_state()}}

As its name suggests, function \texttt{make\_equilibrium\_init\_state()} creates
an equilibrium initial state with a given (possibly zero) temperature $T$ (Sec. \ref{sec:formulation:initial}, points 2 and 3). The function has the following arguments.
\begin{itemize}
    \item $\hat H_0$, Hamiltonian of the system at $t=0$ as an \texttt{Operator} object.
    \item Index set for fermions.
    \item Index set for bosons.
    \item Temperature $T$.
    \item Dictionary of calculation parameters (will be explained shortly).
    \item \texttt{bits\_per\_boson}, number of bits per bosonic
    degree of freedom.
\end{itemize}

In order to construct $\hat\rho_0$, one needs to (partially) diagonalize $\hat H_0$ and obtain some of its eigenpairs $(E_n,\ket{\psi_n})$ with the smallest energies. The diagonalization is preceded by a preparatory step, where $\mathcal{H}$
is partitioned into a direct sum of invariant subspaces of $\hat H_0$,
\begin{equation}
    \mathcal{H} = \mathcal{H}_1 \oplus \mathcal{H}_2 \oplus \ldots \oplus
        \mathcal{H}_S,\quad \hat H_0\mathcal{H}_s = \mathcal{H}_s.
\end{equation}
This sort of space partition effectively transforms $\hat H_0$ to a block-diagonal form and greatly reduces computational costs of the diagonalization. \realevol~employs the first phase of so called autopartition
algorithm described in Sec. 4.2 of \cite{CTHYB}.

Diagonalization of each diagonal block is then performed with one of three
methods, depending on the dimension of the corresponding $\mathcal{H}_s$.
\begin{itemize}
    \item $\dim(\mathcal{H}_s) = 1$, diagonalization of a $1\times1$ matrix is trivial.
    \item $\dim(\mathcal{H}_s)$ is smaller than the value of calculation parameter \texttt{arpack\_min\_matrix\_size}.

    LAPACK routine \texttt{zheev} is called to obtain all eigenpairs of the diagonal block.

    \item $\dim(\mathcal{H}_s)$ is no lesser than \texttt{arpack\_min\_matrix\_size}.

    ARPACK is used to compute a few states with the lowest energies within
    the block. Calculation parameter \texttt{arpack\_tolerance} can be used
    to control the eigenvalue convergence tolerance. Another relevant parameter
    is \texttt{arpack\_ncv}. It is a dictionary \texttt{\{int:int\}} that allows to set NCV (number of Lanczos vectors) for individual subspaces (keys of the dictionary are subspace indices $s$).
\end{itemize}

\texttt{make\_equilibrium\_init\_state()} fills \texttt{InitState} with
eigenvectors of $\hat H_0$ and their statistical weights,
\begin{equation}
    \mathtt{InitState} = \{\mathcal{H}_s\},
    \left\{\left(\ket{\psi_n}, \frac{\exp(-E_n/T)}{Z}, s\right)\right\}.
\end{equation}

The total number of included eigenstates depends on the cutoff imposed on
the statistical weight $\exp(-E_n/T)/Z$. This cutoff can be adjusted via
\texttt{min\_rel\_weight} parameter and equals the machine epsilon by default.
It is worth noting that usually only a few subspaces contribute to the
initial state if the temperature is not too high.

The amount of information being printed out by \texttt{make\_equilibrium\_init\_state()} is controlled via \texttt{verbosity}
parameter.

\bibliographystyle{plain}

\begin{thebibliography}{9}

\bibitem{NoneqDMFT}
``Nonequilibrium dynamical mean-field theory and its applications''.
H. Aoki, N. Tsuji, M. Eckstein, M. Kollar, T. Oka, and Ph. Werner,
Rev. Mod. Phys. {\bf 86}, 779 (2014)

\bibitem{CTHYB}
``TRIQS/CTHYB: A continuous-time quantum Monte Carlo hybridisation
expansion solver for quantum impurity problems``.
P. Seth, I. Krivenko, M. Ferrero and O. Parcollet,
Comp. Phys. Comm. {\bf 200}, 274-284 (2016)

\end{thebibliography}

\end{document}
